\ifdefined\iscombined
	\lecturetitle{Modelling Software Systems}
\else
\documentclass{article}
\usepackage{listings}
\usepackage{amsthm}
\usepackage{comment}
\usepackage{amsmath}
\usepackage{amsfonts}
\usepackage{sectsty}
\usepackage{hyperref}
\usepackage{fancyhdr}
\usepackage[top=1in,bottom=1in,left=1in,right=1in]{geometry}
\usepackage{float}
\usepackage{graphicx}
\usepackage{tabularx}
\usepackage{mathtools}

\date{
	\vspace{-.75em}
	{\large \today}
}

\title{
	\vspace*{-1.5em}
	{\Huge Lecture 14} \\
	\vspace{-1em}
	\rule{\textwidth/4}{.01em} \\
	\vspace{-.25em}
	{\Large Modelling Software Systems}
	\vspace{-.75em}
}

\author{
	{\large Matthew Farah}
}

\hypersetup{
	colorlinks=true,
	linkcolor=blue,
	filecolor=magenta,
	urlcolor=blue,
	citecolor=blue,
	anchorcolor=blue
}

\fancyhead{}
\fancyhead[L]{\today}
\fancyhead[R]{Lecture 14 - Modelling Software Systems}
\sectionfont{\fontsize{12}{12}\bfseries}
\subsectionfont{\fontsize{10}{12}\normalfont}
\subsubsectionfont{\fontsize{10}{12}\normalfont}

\begin{document}

\maketitle
\pagestyle{fancy}

\fi

\begin{itemize}
	\item Models of a software system's design are:
	\begin{enumerate}
		%TODO: Reword
		\item Abstract representations of knowledge of the system.
		\item An approximation of some aspects of the system.
		\item Can be expressed using graphs, natural language, math, and other ways.
	\end{enumerate}
	\item A system's design models are used to:
	\begin{enumerate}
		\item Communicate knowledge about the system.
		\item Help articulate and assess design decisions for the system by abstracting away certain details.
		\item Act as documentation to aid in maintaining the system.
	\end{enumerate}
	\item Design models can be `formal' or `rigorous'.
	\begin{itemize}
		\item Don't let the word rigorous confuse you! Rigorous design models are subject to interpretability.
	\end{itemize}
	\item Examples of formal design models include:
	\begin{enumerate}
		\item FSMs.
		\item Graphs.
		\item Relations.
	\end{enumerate}
	\item Examples of rigorous design include:
	\begin{enumerate}
		\item UML diagrams.
		\item Proprietary/non-standard models.
	\end{enumerate}
	\item Analysis class diagrams are an example of a rigorous design model.
	\item Analysis class diagrams are composed of three types of classes:
	\begin{enumerate}
		\item Boundary classes.
		\begin{itemize}
			\item Are interfaces between the system and its environment
			\item Keep the format of communication between the system and its environment as its secret.
		\end{itemize}
		\item Entity classes.
		\begin{itemize}
			\item Store data.
			\item Exchange data between classes.
			\item Keep an $n$-tuple (representing the data the entity class stores) as its secret.
		\end{itemize}
		%TODO: Elaborate
		\item Control classes.
		\begin{itemize}
			\item Coordinate interactions between other classes in response to a business event.
		\end{itemize}
	\end{enumerate}
	\item The $4 + 1$ model is another example of design modelling.
	\begin{itemize}
		\item The scenario view of the system acts as the basis for its logical view.
		\begin{itemize}
			\item The logical view represents the system as an analysis class diagram.
			\begin{itemize}
				\item All other views are derived from the logical view.
			\end{itemize}
		\end{itemize}
	\end{itemize}
	\item All software architectures can be said to be composed of the following parts:
	\begin{enumerate}
		\item Elements.
		\item Connectors.
		\item The rationale behind those elements/connectors.
	\end{enumerate}
	\item Connectors always have some source and destination element.
	\item `Multiplicity' is an attribute of connectors.
	\begin{itemize}
		\item The multiplicity of a connector describes how many elements are related together by the connector from source to destination.
		\begin{itemize}
			\item The multiplicity of any connector is one of:
			\begin{enumerate}
				\item 1-1.
				\item 1-many.
				\item Many-1.
				\item Many-many.
			\end{enumerate}
		\end{itemize}
	\end{itemize}
	\item `Distance' is another attribute of a connector.
	\begin{itemize}
		\item The distance of a connector describes how `far apart' the elements it connects are.
		\begin{itemize}
			\item How far apart elements are can be categorized as:
			\begin{enumerate}
				\item Elements belonging to the same thread.
				\item Elements belonging to the same process.
				\item Elements belonging to the same computer.
				\item Etc.
			\end{enumerate}
		\end{itemize}
	\end{itemize}
\end{itemize}

\ifdefined\iscombined
	\newpage
\else
	\end{document}
\fi
